Diophantine systems of equations and inequalities arise frequently in Combinatorics.
In Table~\ref{tab:diophant}, Orbiter commands for creating and solving diophantine systems are shown.
\orbitertableofcommands{diophant_create}
In Table~\ref{tab:diophant:activities}, Orbiter activities for diophantine systems are shown.
\orbitertableofcommands{diophant_activity}


Consider the matrix 
$$
A= 
\left[
\begin{array}{cccccc}
0&1&0&1&0&0\\
0&0&1&0&1&0\\
1&0&1&0&0&0\\
0&1&0&1&0&1\\
1&0&0&0&0&1\\
1&0&1&0&0&0\\
0&1&0&0&1&1\\
\end{array}
\right]
$$
Suppose we want to find all column vectors $\bx$ with entries in $0,1$ such that 
$$
A \bx = \bone
$$
where $\bone$ is the all-one column vector.
Orbiter offers two algorithms to do this. 
One is McKay's possolve, the other is Knuth's dancing links (DLX).
In order to get started, we need to create a diophant object.
In the following example, we use the makefile variable \verb'TEST_SYSTEM' for the 
coefficient matrix. 

\sourcecode{TEST_SYSTEM}

The right hand side will be specified using the \verb'RHS' command below:

\sourcecode{solve_test_system}


There are two commands to solve a diophantine system: \verb'-solve_mckay' and \verb'-solve_DLX'. The latter is more restrictive, 
as it allows only 0,1 variables.
Here is the McKay solver:

\sourcecode{McKay_test}

The solutions are written to the file \verb'DLX_test.sol'. 
And now the dancing links solver:

\sourcecode{DLX_test}



\bigskip

